\par\medskip\noindent\textbf{\task{conditional\_shift}, letters ($N = 8$, $k = 4$, $A = 4$).}\par
\begin{small}
\begin{verbatim}
You are a deterministic sequence transformer.
Task: Branching character shift.
Alphabet: A, B, C, D (positions numbered 0, 1, 2, ... in this order)
Rule: For each character, look at its position number in the alphabet.
- If the position is EVEN, replace the character with the NEXT character in the alphabet
(e.g. A -> B).
- If the position is ODD, replace the character with the PREVIOUS character in the alphabet
(e.g. B -> A).
Wrap around the alphabet if needed.
Constraint: The input contains vertical bars (|). These are structural boundaries. DO NOT
shift them. Copy them exactly to the output.
No spaces, no explanations.

Input: DADB|DAAD
Output:            (expected: CBCA|CBBC)
\end{verbatim}
\end{small}

\par\medskip\noindent\textbf{\task{conditional\_shift}, digits: lines that differ from the letter prompt ($N = 8$, $k = 4$, $A = 4$).}\par
\begin{small}
\begin{verbatim}
Alphabet: 0, 1, 2, 3 (positions numbered 0, 1, 2, ... in this order)
- If the position is EVEN, replace the character with the NEXT character in the alphabet
(e.g. 0 -> 1).
- If the position is ODD, replace the character with the PREVIOUS character in the alphabet
(e.g. 1 -> 0).
Input: 3031|3003
\end{verbatim}
\end{small}

\par\medskip\noindent\textbf{\task{random\_permutation}, letters ($N = 8$, $k = 4$, $A = 4$).}\par
\begin{small}
\begin{verbatim}
You are a deterministic sequence transformer.
Task: Character substitution mapping.
Alphabet: A, B, C, D
Mapping: A->D, B->C, C->A, D->B
Rule: Replace each character with its mapped character. This is a fixed one-to-one
substitution cipher.
Constraint: The input contains vertical bars (|). These are structural boundaries.
DO NOT transform them. Copy them exactly to the output.
No spaces, no explanations.

Input: DADB|DAAD
Output:            (expected: BDBC|BDDB)
\end{verbatim}
\end{small}

\par\medskip\noindent\textbf{\task{running\_xor}, original wording ($N = 8$, $k = 4$).}\par
\begin{small}
\begin{verbatim}
You are a deterministic sequence transformer.
Task: Running XOR with previous input.
Alphabet: 0, 1 (binary digits)
Rule: At each position, output the XOR of the current input digit with the previous input
digit.
The first digit is copied as-is (no previous input to XOR with).
Each subsequent digit: output[i] = input[i] XOR input[i-1]
Example: Input 1011 -> Output 1110 (calculation: 1, 0^1=1, 1^0=1, 1^1=0)
Constraint: The input contains vertical bars (|). These are structural boundaries.
DO NOT transform them. Copy them exactly to the output.
No spaces, no explanations.

Input: 1001|1111
Output:            (expected: 1101|0000)
\end{verbatim}
\end{small}

\par\medskip\noindent\textbf{\task{running\_xor}, reset variant: lines that differ ($N = 8$, $k = 4$).}\par
\begin{small}
\begin{verbatim}
Task: Running XOR with previous input, reset at boundaries.
IMPORTANT: The input contains vertical bars (|), structural boundaries. Immediately after
each boundary, treat that position exactly like the first digit of the whole input -- copy
it as-is (nothing before a boundary counts as "previous"), then resume the XOR rule from
there until the next boundary.
Example: Input 1011|0010 -> Output 1110|0011 (first block: 1 copied, 0^1=1, 1^0=1, 1^1=0;
second block restarts: 0 copied, 0^0=0, 1^0=1, 0^1=1)
Copy the boundary characters themselves exactly to the output.
\end{verbatim}
\end{small}

\par\medskip\noindent\textbf{\task{running\_xor}, neutral wording: lines that differ ($N = 8$, $k = 4$).}\par
\begin{small}
\begin{verbatim}
Constraint: The input contains vertical bars (|). Copy each bar to the output at the same
position.
\end{verbatim}
\end{small}

\par\medskip\noindent\textbf{\task{running\_xor}, across wording: lines that differ ($N = 8$, $k = 4$).}\par
\begin{small}
\begin{verbatim}
Constraint: The input contains vertical bars (|). They are only visual separators: apply the
rule to the digits as if the bars were not there (the digit right after a bar uses the digit
right before the bar as its previous input), and copy each bar to the output at the same
position.
\end{verbatim}
\end{small}

\par\medskip\noindent\textbf{\task{running\_xor}, across wording with worked examples: lines that differ ($N = 8$, $k = 4$).}\par
\begin{small}
\begin{verbatim}
Constraint: The input contains vertical bars (|). They are only visual separators: apply the
rule to the digits as if the bars were not there, and copy each bar to the output at the
same position. The digit right after a bar still depends on the input digit(s) right before
the bar.
Worked examples with bars:
Input 01|01|00|11 -> Output 01|11|10|10 (the output digit right after the first bar is 1: it
uses the input digit just before the bar (1) together with the digit right after it (0))
Input 0101|1000 -> Output 0111|0100
\end{verbatim}
\end{small}

\par\medskip\noindent\textbf{\task{copy\_previous}: lines that differ from \task{running\_xor} ($N = 8$, $k = 4$).}\par
\begin{small}
\begin{verbatim}
Task: Copy the previous input digit.
Rule: At each position, output the PREVIOUS input digit (not the current one).
The first digit is copied as-is (no previous input to copy).
Each subsequent digit: output[i] = input[i-1]
Example: Input 1011 -> Output 1101 (calculation: 1 (copied), then 1, 0, 1 -- each shifted
from the position before)
\end{verbatim}
\end{small}

\par\medskip\noindent\textbf{\task{running\_sum\_local}: lines that differ from \task{running\_xor} ($N = 8$, $k = 4$, $M = 4$).}\par
\begin{small}
\begin{verbatim}
Task: Running sum with previous input (local window).
Alphabet: 0-3 (digits)
Modulus: 4
Rule: At each position, output the sum of the current input digit and the previous input
digit, modulo 4.
The first digit is copied as-is (no previous input to add).
Each subsequent digit: output[i] = (input[i] + input[i-1]) mod 4
Input: 2122|0101
\end{verbatim}
\end{small}

\par\medskip\noindent\textbf{\task{sliding\_window\_max}: lines that differ from \task{running\_xor} ($N = 8$, $k = 4$, $A = 4$).}\par
\begin{small}
\begin{verbatim}
Task: Sliding-window maximum.
Alphabet: 0-3
Rule: At each position, output the LARGEST digit among this position and the 2 position(s)
immediately before it (a window of the last 3 digits, including the current one).
Near the very start of the sequence, where fewer than 3 digits exist yet, use only the
digits that exist so far.
Example (window=3): Input 3021 -> Output 3332 (position 0: window=[3], max=3. position 1:
window=[3,0], max=3. position 2: window=[3,0,2], max=3. position 3: window=[0,2,1] (oldest
'3' has fallen out of the window), max=2.)
DO NOT include them in the window or count them as digits. Copy them exactly to the output.
Input: 2122|0101
\end{verbatim}
\end{small}

\par\medskip\noindent\textbf{\task{dna\_copy\_previous} (real promoter DNA) ($N = 8$, $k = 4$).}\par
\begin{small}
\begin{verbatim}
You are a deterministic sequence transformer.
Task: Copy the previous base.
Alphabet: A, C, G, T in the input (a stretch of a human promoter sequence); A, C, G, T and N
in the output.
Rule: At each position, output the base that comes immediately before it in the input. The
first position has no previous base and gets N.
Example: Input ACGTTG -> Output NACGTT
Constraint: The input contains vertical bars (|). These are structural boundaries.
DO NOT transform them. Copy them exactly to the output.
No spaces, no explanations.

Input: CTGG|CGGA
Output:            (expected: NCTG|GCGG)
\end{verbatim}
\end{small}

\par\medskip\noindent\textbf{\task{dna\_same\_class}: lines that differ from \task{dna\_copy\_previous} ($N = 8$, $k = 4$).}\par
\begin{small}
\begin{verbatim}
Task: Same base class as the previous base.
Alphabet: A, C, G, T in the input (a stretch of a human promoter sequence); digits 0 and 1
in the output.
Rule: A and G are purines; C and T are pyrimidines. For each base, output 1 if it is in the
same class as the base immediately before it, and 0 otherwise. The first base has no
previous base and gets 0.
Example: Input AGCTTA -> Output 010110 (A: first, 0; G: purine after purine, 1; C: pyrimidine
after purine, 0; T: pyrimidine after pyrimidine, 1; T: pyrimidine after pyrimidine, 1; A:
purine after pyrimidine, 0)
\end{verbatim}
\end{small}

\par\medskip\noindent\textbf{\task{text\_copy\_previous} (WikiText-2 letters) ($N = 8$, $k = 4$).}\par
\begin{small}
\begin{verbatim}
You are a deterministic sequence transformer.
Task: Copy the previous letter.
Alphabet: lowercase letters a-z in the input (English text with spaces and punctuation
removed); the same letters and # in the output.
Rule: At each position, output the letter that comes immediately before it in the input. The
first position has no previous letter and gets #.
Example: Input thelazydog -> Output #thelazydo
Constraint: The input contains vertical bars (|). These are structural boundaries.
DO NOT transform them. Copy them exactly to the output.
No spaces, no explanations.

Input: toth|ealr
Output:            (expected: #tot|heal)
\end{verbatim}
\end{small}

\par\medskip\noindent\textbf{\task{text\_ascending}: lines that differ from \task{text\_copy\_previous} ($N = 8$, $k = 4$).}\par
\begin{small}
\begin{verbatim}
Task: Later in the alphabet than the previous letter.
Alphabet: lowercase letters a-z in the input (English text with spaces and punctuation
removed); digits 0 and 1 in the output.
Rule: For each letter, output 1 if it comes later in the alphabet than the letter
immediately before it, and 0 otherwise (also 0 if the two letters are the same). The first
letter has no previous letter and gets 0.
Example: Input thelazydog -> Output 0001010010 (t: first, 0; h after t: 0; e after h: 0; l
after e: 1; a after l: 0; z after a: 1; y after z: 0; d after y: 0; o after d: 1; g after o:
0)
\end{verbatim}
\end{small}
